\documentclass[12pt]{article}
\usepackage{graphicx} % Required for inserting images
\usepackage[utf8]{inputenc} % Codificación de caracteres
\usepackage[spanish]{babel} % Idioma del documento
\usepackage{amsmath, amssymb, amsthm}
\usepackage{pgfplots}
\usepackage{tabularx}
\usepackage[margin=1in]{geometry}
\usepackage{svg}
\usepackage{float}
\usepackage{fancyhdr}
\usepackage{setspace}
\usepackage[most]{tcolorbox}

\newtcolorbox{tadelis}{
    colback=blue!5,
    colframe=blue!60!black,
    title=\textit{Game Theory, Steven Tadelis},
    fonttitle=\bfseries,
    boxrule=0.8pt,
    arc=2mm
}

\newtcolorbox{comentarioprofe}{
    colback=green!5,
    colframe=green!50!black,
    title=Comentario,
    fonttitle=\bfseries,
    boxrule=0.8pt,
    arc=2mm
}

\newtcolorbox{importante}{
    colback=orange!8,
    colframe=orange!70!black,
    title=Importante,
    fonttitle=\bfseries,
    boxrule=0.8pt,
    arc=2mm
}


\geometry{margin=1in}
\onehalfspacing
\pagestyle{fancy}
\rhead{Teoría de Juegos e Información}
\lhead{Incertidumbre y Decisiones Racionales}
\title{Tema 1. Incertidumbre y Decisiones Racionales}
\author{Ignacio Sanabria}
\date{Agosto 26}
\begin{document}
\maketitle

\section{¿Que es la teoría de juegos?}
La teoria de juegos busca explicar la forma en que se toman decisiones racionales en el día a día. \\
\textbf{Agente Racional:} Un agente racional es aquel que busca siempre obtener su mayor utilidad. \\
\subsection{Conceptos Básicos}
\begin{itemize}
    \item Triviales: no importan tanto, son superficiales.
    \item Trascendentales: Afectan mis pagos esperados a futuro. 
    \item \textbf{Supuestos:} vitales para explicar un modelo. Un exceso de supuestos convierte el modelo en inexacto. 
\end{itemize}

\section{Incertidumbre}
Se platea un problema de decisión que incluye: 
\begin{itemize}
    \item Situaciones sobre las que se debe decidir. 
    \item Conjunto de agentes económicos: Jugadores $J=\{\}$
    \item Conjunto de acciones a tomar: Acciones $A=\{\}$
    \item Conjunto de resultados: Resultados $R=\{\}$
\end{itemize}
Para que un problema de decisión tenga sentido, los agentes económicos deben tener preferencias sobre los resultados. \\
Estas vienen dadas por: 
\begin{enumerate}
    \item $x\succeq y \xrightarrow{}$ Débilmente Preferido 
    \item $x\succ y \xrightarrow{}$ Estrictamente preferido
    \item $x\sim y \xrightarrow{}$ Indiferente
\end{enumerate}

\section{Axiomas de Preferencias}
\begin{enumerate}
    \item \textbf{Axioma de Completitud:} Los agentes económicos siempre son capaces de ordenar sobre dos resultados. \texttit{No pueden quedar indiferentes ante dos resultados.}
    \item \textbf{Axioma de Transitividad:} Son capaces de ordenar todos los resultados. \textit{Garantiza que no existen contradicción de elecciones de los jugadores.} 
\end{enumerate}
\subsection{Paradoja de Condorcet 1785}
Condorcet planteo que es posible al tener un grupo de agentes individualmente racionales, y ponerlos a tomar una decision en grupo, estos empiecen a tomar decisiones irracionales como grupo. \\
Sea $J=\{A,B,C\}$ un conjunto de estudiantes que deben elegir si prefieren un quiz teorico ($t$), practico ($p$), o mixto ($m$). 
\begin{align*}
    \text{Con preferencias:}&\\\\
    A:& \quad t\succ m\succ p \\
    B: & \quad m \succ p \succ t \\
    C: & \quad p \succ t \succ m 
\end{align*}
Si les pedimos votar obtendriamos: 
\begin{align*}
    \text{t vs m:}& \quad t \succ m \quad \textit{A favor de t: A,C}\\
    \text{p vs m:}& \quad m \succ p \quad \textit{A favor de m: A,B}\\
    \text{t vs p:}& \quad p \succ t \quad \textit{A favor de p: B,C}\\
\end{align*}
Sin embargo, esto rompe el \textbf{axioma de transitividad} dado que:
\begin{align*}
   \text{Por transitividad:}& \quad  t \succ m \succ p  \\
   \text{pero la ultima votacion coloca:}&\quad  p \succ t \\
   \therefore t \succ p \;\mathbf{\Rightarrow \Leftarrow} \;p \succ t
\end{align*}

\begin{tadelis}
Según el libro, la ecuación de Euler caracteriza
la senda óptima de consumo.
\end{tadelis}
\begin{comentarioprofe}
El profesor enfatizó que esta condición solamente se cumple
bajo determinadas condiciones de interioridad.
\end{comentarioprofe}
\begin{importante}
La tasa de interés afecta el costo de oportunidad
del consumo presente.
\end{importante}

\end{document}
